% Aionis — arXiv preprint scaffold (PREP, not yet submitted).
%
% Framing (a): the power-floor theorem is the lead methodological contribution,
% supported by the end-to-end anti-leakage discipline and the 15-row null evidence.
%
% Build: latexmk -pdf main.tex   (or upload this folder to Overleaf)
% No local TeX toolchain in the Aionis repo; compile on a machine with TeX Live
% (or Overleaf). Uses ONLY standard, universally-available packages (no custom .cls).
% To target an Elsevier venue later, swap \documentclass{article} -> elsarticle
% (download elsarticle.cls from CTAN/Overleaf); content is class-agnostic.
%
% Status: v0.1-prep 2026-08-05. Content cross-checked against ledger #49 and the
% Chinese v1.0 / English v1.0-en drafts. Awaiting owner authorization before any
% arXiv upload (irreversible externalization).
\documentclass[11pt,a4paper]{article}

\usepackage[margin=1in]{geometry}
\usepackage{amsmath,amssymb}
\usepackage{booktabs}
\usepackage{graphicx}
\usepackage{hyperref}
\usepackage[round,authoryear]{natbib}
\usepackage{xcolor}
\hypersetup{colorlinks=true,linkcolor=blue!60!black,citecolor=blue!60!black,urlcolor=blue!60!black}

\title{A Power Floor for Monthly Cross-Sectional Rank-IC Equivalence Testing:\\
An Anti-Leakage Discipline Case Study on U.S.\ and A-Share Equities}

\author{Rethymus\thanks{Corresponding author. The Aionis project is MIT-licensed; the
research harness, frozen configurations, and append-only ledger are available at
\href{https://github.com/Rethymus/Aionis}{github.com/Rethymus/Aionis}. Data:
see \S\ref{sec:replication}.}}\date{Preprint v0.1 \today}

\begin{document}
\maketitle

\begin{abstract}
We document a \emph{power floor} for monthly cross-sectional rank-IC equivalence
testing: under a pre-registered Jennison--Turnbull group-sequential equivalence
gate at SESOI $\pm 0.010$, the empirically documented noise floor of monthly
rank-IC ($\sigma\!\approx\!0.10$, consistent with the predictability levels of
\citealt{gu2020empirical} and the Fundamental Law of \citealt{grinold1999active})
makes declaring equivalence structurally infeasible within realistic sample sizes
(look-3 $n_{\min}=435$ months, $\approx 36$ years). We establish this as a
general property---not an artifact of our data---by verifying $\sigma(\text{IC})\in[0.092,0.163]$
across 20 of our own IC series spanning single- and dual-region configurations,
all within the documented $[0.08,0.15]$ range. As a carrier for the theorem we
present Aionis, a falsifiable, anti-leakage research harness whose disciplines
(config-before-result append-only ledger, full point-in-time data, purged
cross-validation with embargo, and bit-identical determinism) are themselves the
object of study. On S\&P 500 point-in-time constituents and A-share CSI 300, 15
pre-registered two-tailed configurations yield null point estimates, including
the first confirmatory out-of-sample run (combined rank-IC $-0.0088$, $p=0.484$).
The pre-registered Jennison--Turnbull gate returns \textsc{Not\_Equivalent} at
look-1 even with a null point estimate---a live demonstration that the discipline
refuses to prematurely declare equivalence. The honest, pre-registered contribution
is therefore the null, the end-to-end discipline, and the power-limit disclosure:
not a declared equivalence.
\end{abstract}

\section{Introduction}\label{sec:intro}

The common failure mode of quantitative asset-pricing research is not ``no alpha
found'' but \emph{unfalsifiability}: results silently rewritten via rerun-to-significance,
backfilling, survivor bias, or cross-validation leakage. This preprint makes two
contributions against that backdrop.

First, we document a \emph{power floor}: pre-registered equivalence testing on the
monthly cross-sectional rank-IC at a $\pm 0.010$ equivalence margin is structurally
infeasible within realistic sample sizes, because the noise floor of monthly rank-IC
($\sigma\!\approx\!0.10$) dominates the margin by an order of magnitude. We ground
this floor in the published predictability literature
\citep{gu2020empirical,goyal2008comprehensive,grinold1999active} and verify it
empirically across 20 of our own IC series (\S\ref{sec:powerfloor}).

Second, we present the discipline that makes the null \emph{credible}: an
append-only frozen-configuration ledger (sha256 committed before any out-of-sample
metric is observed), a full point-in-time data stack, purged cross-validation with
embargo, and bit-identical determinism. The first confirmatory out-of-sample run
returns a null point estimate, and the pre-registered equivalence gate refuses to
declare equivalence at look-1---the discipline operating exactly as designed.

Throughout, \emph{null is the intended, publishable outcome}. We do not claim market
efficiency, do not claim factor alpha is strictly zero, and do not constitute a
tradable strategy (see \S\ref{sec:limitations}).

\section{Related Literature}\label{sec:related}

\citealt{gu2020empirical} establish that even the strongest machine-learning models
achieve monthly out-of-sample $R^2$ of only $1.08$--$1.80\%$ on U.S.\ equities,
implying mean rank-IC magnitudes of order $0.05$--$0.12$ (since $R^2\!\approx\!\text{IC}^2$
for a single cross-sectional signal). \citealt{goyal2008comprehensive} comprehensively
re-examine equity-premium predictors and find that most fail out-of-sample. The
Grinold--Kahn Fundamental Law of Active Management \citep{grinold1999active} frames
the information ratio (IR$\,=$ mean IC$/\sigma(\text{IC})\times\sqrt{\text{breadth}}$);
an annualized IR of $0.5$ is conventionally ``good,'' which implies mean IC$/\sigma(\text{IC})\!\approx\!0.14$
and thus $\sigma(\text{IC})\!\approx\!7\!\times\!\text{mean IC}$. For a typical weak
factor (mean IC$\approx 0.015$), this gives $\sigma(\text{IC})\!\approx\!0.10$---exactly
the order of magnitude we estimate. Equivalence testing via TOST
\citep{schuirmann1987,lakens2017} is methodologically standard, but we are unaware
of a prior finance equivalent-margin convention for rank-IC; our $\pm 0.010$ is a
stringent choice informed by practitioner IC norms.

\section{The Anti-Leakage Discipline}\label{sec:discipline}

Aionis structurally locks the leakage sources that normally undermine falsifiability.
Each item below is auditable (it corresponds to a ledger row, an ADR, or a code assertion),
not narrative.

\begin{enumerate}
\item \textbf{\texttt{config\_committed} before result.} The sha256 of the frozen
configuration is appended to an append-only ledger (\texttt{runs/ledger.jsonl})
\emph{before} any out-of-sample metric is observed. Changing the configuration creates
a new ledger row; it never silently overwrites. A headline cannot be
rerun-to-significance rescued.
\item \textbf{Full point-in-time (PIT) stack.} Fundamentals are aligned by
\texttt{filed} date (not period-end); macro variables use ALFRED as-of vintages;
VIX uses a no-revision contract; S\&P 500 membership uses PIT constituents
($\texttt{constituents\_on}(t)$, never a today-snapshot); A-share prices use raw
adjustment factors (frozen strategy).
\item \textbf{Purged validation with embargo.} Historical phases use
\texttt{PurgedGroupKFold} (group = month, embargo = 21 sessions); the chronological
tracks (B, C) use walk-forward folds with train strictly preceding test.
\item \textbf{$H_6$ determinism.} Single-threaded execution ($n\_jobs=1$), all seeds
pinned to $0$, and version-pinned dependencies (\texttt{uv.lock}) make both the IC
series and the raw out-of-sample scores bit-identical across reruns (asserted).
\item \textbf{Two-tailed, pre-registered, null-favored.} One pre-registered
two-tailed claim per phase; SESOI $\pm 0.010$; a Jennison--Turnbull group-sequential
equivalence gate \citep{jennison2000group} with look-specific repeated confidence
intervals (RCI) at $99.44/97.64/95.00\%$ and strict containment.
\item \textbf{Multiplicity budget $= 1$.} The conditioning in Track C is a single
pre-specified interaction (score $\times$ regime state), structurally bypassing the
``tried $K$ subgroups and reported the most significant.''
\end{enumerate}

\section{The Dual-Region Joint Estimator}\label{sec:estimator}

The estimator throughout is the monthly cross-sectional rank-IC (Spearman correlation
of model scores with realized forward returns), with HAC (Newey--West) inference.
Track C introduces a dual-region joint walk-forward over U.S.\ (S\&P 500) and A-share
(CSI 300) month-end cross-sections, aligned by calendar month.

\paragraph{Core theorem.} Month-end sampling with calendar-fold boundaries implies
that the per-region 21-session embargo is automatically satisfied: adjacent month-ends
in either market are $\approx 21$ sessions apart, so no per-session embargo bookkeeping
is required.

\paragraph{Region-month group.} The ranking query is (region, month): U.S.\ names
rank only against U.S.\ names and A-share names only against A-share names, eliminating
cross-currency label contamination (forward returns are local-currency). The model is
jointly fitted with shared parameters.

\paragraph{Anti-leakage invariants (code-asserted).} Per-region chronology
($\max(\text{train}) < \min(\text{test})$ per region), train-only binner fit,
regime state computed as-of month-end (never retrospectively recomputed),
region-month group separation, and $H_6$ bit-identity.

\section{The Power Floor}\label{sec:powerfloor}

\subsection{Prospective power analysis}

Calibrating from the confirmatory IC series (combined $\sigma\!\approx\!0.106$,
lag-1 $\rho\!\approx\!0.07$, observed $\text{se}_\text{HAC}=0.0126$ at $n=71$), the
Jennison--Turnbull schedule $\{60, 90, 120\}$ with O'Brien--Fleming spending
($z_k = 2.772, 2.263, 1.960$; RCI levels $99.44/97.64/95.00\%$) yields the following
minimum sample sizes to declare equivalence at SESOI $\pm 0.010$:

\begin{center}
\begin{tabular}{lrrr}
\toprule
Look & $z_k$ & $n_{\min}$ (months) & years \\
\midrule
1 ($n=60$) & 2.772 & 869 & 72.5 \\
2 ($n=90$) & 2.263 & 580 & 48.3 \\
3 ($n=120$) & 1.960 & 435 & 36.2 \\
\bottomrule
\end{tabular}
\end{center}

Under a 2000-replication block bootstrap, $P(\text{equivalence}) = 0.0000$ at all
three planned looks; the look-3 RCI half-width median ($0.019$) is twice the SESOI.
Declaring equivalence is therefore unreachable within any realistic study horizon.

\subsection{Literature anchoring and cross-configuration empirical verification}

The $\sigma\!\approx\!0.10$ floor is not an artifact of the confirmatory run.
It is consistent with the monthly predictability levels of
\citet{gu2020empirical} ($R^2 \in [1.08\%, 1.80\%]$) and the Fundamental Law
inference from \citet{grinold1999active}. We further verify it directly across 20
of our own IC series---Track C (4 dual-region configurations $\times$ \{U.S., CN,
combined\}) and Track B / Phase B (4 independent run signatures $\times$
\{treatment, baseline\})---finding $\sigma(\text{IC}) \in [0.092, 0.163]$ with median
$0.109$, all within the documented $[0.08, 0.15]$ range. The confirmatory combined
$\sigma = 0.106$ sits at the median. The power floor is thus both literature-anchored
and empirically replicated across our configurations.

\subsection{Theoretical pure-noise bound and the ML noise excess}

The cross-sectional Spearman rank-IC under the null of no predictability has the
closed-form standard deviation $\sigma_{\text{null}} = 1/\sqrt{N_{\text{cross}}-1}$
for a cross-section of $N_{\text{cross}}$ assets. For our cross-sections
($N \in \{462, 465, 929, 1386\}$), this gives $\sigma_{\text{null}} \in [0.027, 0.047]$---yet
the observed $\sigma(\text{IC})$ across 21 series is consistently
\emph{2.0--4.4$\times$} this bound (median excess $3.4\times$; minimum $2.03\times$).
The power floor is therefore not a pure statistical necessity: at the pure-noise
bound $\sigma \approx 0.047$, look-3 would be marginally powered ($n_{\min} \approx 83$
months $< 120$); the floor binds because real ML-based rank-IC carries a robust
noise excess (model fitting noise, cross-sectional heteroskedasticity, and
overlapping forward returns). This refines the claim from ``$\sigma \approx 0.10$ is
a floor'' to a mechanistic statement: the floor is the gap between the closed-form
sampling bound and the empirically robust ML noise excess.

\section{Evidence}\label{sec:evidence}

Table~\ref{tab:evidence} reports 15 pre-registered two-tailed configurations on the
S\&P 500 PIT universe (588 clean tickers, 2016+) and the dual-region Track C panel.
\emph{Honest grading}: shared-fold purged cross-fitting is a CV-proxy, not
chronological out-of-sample; exploratory is not confirmatory. The Grade column is the
true upper bound of publishable strength.

\begin{table}[h]
\centering
\small
\caption{15 null configurations. All CIs bracket zero.}
\label{tab:evidence}
\begin{tabular}{rlrrrrl}
\toprule
\# & Result & Estimate & 95\% CI & $p$ & $n$ & Grade \\
\midrule
1  & Phase B differential        & $-0.0008$ & $[-0.0106, 0.0090]$ & 0.87 & 125 & CV-proxy \\
2  & Phase C differential        & $-0.0065$ & $[-0.0195, 0.0066]$ & 0.36 & 125 & CV-proxy \\
3  & Phase D differential        & $-0.0030$ & $[-0.0137, 0.0078]$ & 0.60 & 125 & CV-proxy \\
4  & Phase E1 differential       & $-0.0028$ & $[-0.0115, 0.0059]$ & 0.53 & 125 & CV-proxy \\
5  & Track B treatment (\#41)    & $+0.0055$ & $[-0.021, 0.033]$   & 0.69 & 125 & chron./explor. \\
6  & Track B differential        & $+0.0076$ & $[-0.004, 0.020]$   & 0.22 & 125 & chron./explor. \\
7  & Track C joint combined IC   & $-0.0070$ & $[-0.030, 0.016]$   & 0.55 &  71 & chron./explor. \\
8  & Track C conditional-IC $\beta$ & $-0.015$ & ---               & 0.20 &  71 & explor. \\
9  & Track C 3-layer cond.-IC    & $-0.0148$ & ---                 & 0.21 &  71 & explor. \\
10 & Baseline-FF5                & $+0.0106$ & half $0.0196$       & ---  & 125 & CV-proxy \\
11 & Baseline-RANK               & $+0.0154$ & half $0.0149$       & 0.04 & 125 & CV-proxy \\
12 & $h\in\{10,42\}$ sweep (8)   & all bracket zero & ---          & 0.41--0.86 & 124--126 & explor. \\
13 & Track C asymmetric-35       & $-0.0121$ & $[-0.035, 0.011]$   & 0.31 &  71 & explor. \\
14 & Track C asymmetric-41       & $-0.0088$ & $[-0.034, 0.016]$   & 0.48 &  71 & explor. \\
\textbf{15} & \textbf{Track C confirmatory} & $\mathbf{-0.0088}$ & $[-0.034, 0.016]$ & \textbf{0.48} & \textbf{71} & \textbf{confirmatory} \\
\bottomrule
\end{tabular}
\end{table}

\section{First Confirmatory Out-of-Sample}\label{sec:confirmatory}

The first confirmatory OOS run (ledger row \#49, frozen configuration \#48 committed
before observation) yields:
\begin{itemize}
\item Combined rank-IC mean $-0.00884$ ($p_\text{HAC}=0.484$; 95\% HAC CI $[-0.034, 0.016]$
brackets zero; $n=71$ months).
\item Per-region: U.S.\ IC $+0.0052$, CN IC $-0.0265$ (both null).
\item Conditional-IC $\beta = -0.0076$ ($p = 0.43$): no regime interaction.
\item Jennison--Turnbull look-1 ($n=60$, RCI $99.44\%$): \textsc{Not\_Equivalent};
RCI $[-0.051, +0.027]$ is wider than $\pm 0.010$ SESOI.
\item $H_6$ double-run bit-identical: \textsc{Pass}.
\end{itemize}

The honest reading is a null point estimate plus an underpowered look-1---neither
``treatment works'' nor ``equivalence rejected.'' The gate's refusal to prematurely
declare equivalence, even with a null point estimate, is the live demonstration of
the discipline. Strict equivalence (look-2/3) requires E3 forward-live calendar time
(look-2 $\approx 2028$, look-3 $\approx 2031$).

\section{Limitations}\label{sec:limitations}

\begin{itemize}
\item CV-proxy (shared-fold purged cross-fit) is not chronological OOS; purge/embargo
prevent label overlap but the train complement may include post-test months.
\item Exploratory is not confirmatory; only row 15 is confirmatory.
\item The power floor (\S\ref{sec:powerfloor}) means strict equivalence at $\pm 0.010$
is unreachable within a realistic horizon; the contribution is the null, the discipline,
and the power-limit disclosure---not a declared equivalence.
\item PIT constituents mitigate survivor bias but cannot eliminate it (no free
delisting PIT). Currency: the region-month group eliminates label cross-contamination,
but the joint IC is still a composition of local-currency returns.
\item Not investment advice. Null does not prove market efficiency, does not prove
factor alpha is strictly zero, and does not constitute a tradable strategy.
\end{itemize}

\section{Reproducibility}\label{sec:replication}

Aionis is reproducible by construction. The append-only ledger records every frozen
configuration's sha256 before any result; dependencies are version-pinned
(\texttt{uv.lock}); all seeds are $0$; execution is single-threaded. Independent
parties clone the MIT-licensed repository, supply their own API keys (Tiingo, FRED),
run the tracked fetch scripts to rebuild the point-in-time cache, and reproduce any
frozen run bit-identically (the $H_6$ assertion fails on any drift). Raw vendor data
are not redistributed (vendor terms); public-domain sources (SEC EDGAR, FRED/ALFRED)
are redistributable. Full statement: \texttt{docs/replication-availability.md}.

\section{Conclusion}\label{sec:conclusion}

We document that monthly cross-sectional rank-IC equivalence testing at $\pm 0.010$
hits a structural power floor ($\sigma\!\approx\!0.10$ noise vs.\ $\pm 0.010$ margin,
requiring $\approx 36$ years at the final look), establish it as a general property
via literature anchoring and 20-series empirical verification, and present a
confirmatory-grade null carried end-to-end by an anti-leakage discipline whose gate
refuses to prematurely declare equivalence. The contribution is the null, the
discipline, and the disclosed power limit.

\bibliographystyle{plainnat}
\bibliography{references}

\end{document}
